% chapters/chapter5_results.tex
\chapter{Results}
\label{ch:results}

% ============================================================
% CHAPTER OVERVIEW
% ============================================================
% Present results organized by research question:
% - RQ1: What neural/artifactual signatures are present?
% - RQ2: What features does EEGNet learn?
% - RQ3: What drives classification decisions?
%
% End with hypothesis scoring and correlation analysis
%
% TARGET: 15-20 pages total
%
% STYLE NOTES:
% - Present data objectively without interpretation
% - Save interpretation for Discussion chapter
% - Use figures and tables extensively
% - State statistical results precisely (p-values, effect sizes)


% ============================================================
% 5.1 NEUROPHYSIOLOGICAL CHARACTERIZATION (RQ1)
% ============================================================
\section{Neurophysiological Characterization of the Dataset}
\label{sec:results_neuro}

% TARGET: 5-6 pages


% ------------------------------------------------------------
\subsection{Grand Average ERDS}
\label{subsec:results_grand_erds}

% TARGET: 2-2.5 pages
%
% CONTENT:
%
% FIGURE: Grand average time-frequency ERDS plot
% - Channels: C3 and C4 (side by side)
% - Time: -2s to 3s relative to cue
% - Frequency: 1-40 Hz
% - Color: ERDS percentage (blue = ERD, red = ERS)
% - Mark cue onset (t=0) with vertical line
%
% OBSERVATIONS TO REPORT:
%
% Mu band (8-13 Hz):
% - Is ERD present? At what channels?
% - When does ERD begin? (Before or after cue?)
% - Is ERD sustained through imagery period?
% - Is there contralateral lateralization? (C3 for right MI, C4 for left)
%
% Beta band (13-30 Hz):
% - Is ERD present?
% - Is there post-movement beta rebound (ERS)?
%
% Low frequency (<4 Hz):
% - Any activity? Could indicate MRCP or artifacts
%
% Pre-cue vs Post-cue:
% - Compare ERDS magnitude in [-1, 0]s vs [0, 2]s
% - Does pre-cue show discriminative activity?
%
% STATISTICAL RESULTS:
% - Cluster-based permutation test: left vs right MI
% - Report significant time-frequency clusters
% - Report effect sizes
%
% INCLUDE FIGURE:
% - 2x2 panel: Left MI/C3, Left MI/C4, Right MI/C3, Right MI/C4
% - Or: Difference map (Left - Right) at C3 and C4


% ------------------------------------------------------------
\subsection{Topographic Distribution}
\label{subsec:results_topography}

% TARGET: 1.5-2 pages
%
% CONTENT:
%
% FIGURE: Topographic maps of ERDS
% - Multiple panels at key time-frequency points:
%   * Pre-cue (-1.0 to -0.5s), Mu band
%   * Post-cue (0.5 to 1.0s), Mu band
%   * Post-cue (0.5 to 1.0s), Beta band
% - Separate maps for Left MI and Right MI
%
% OBSERVATIONS TO REPORT:
%
% Spatial distribution:
% - Is ERD localized to sensorimotor cortex (C3/C4)?
% - Is there contralateral dominance?
% - Is there any frontal activity? (artifact concern)
% - Is there any occipital activity? (visual concern)
%
% Pre-cue vs Post-cue:
% - Does the topography differ?
% - Is pre-cue activity sensorimotor or elsewhere?
%
% INCLUDE FIGURE:
% - Grid of topographic maps
% - Clear labeling of time window and frequency band


% ------------------------------------------------------------
\subsection{MRCP Analysis}
\label{subsec:results_mrcp}

% TARGET: 1-1.5 pages
%
% CONTENT:
%
% FIGURE: MRCP waveforms
% - Channels: Cz (central), C3, C4 (lateral)
% - Time: -2s to 2s relative to cue
% - Show mean ± SEM across subjects
% - Separate for Left and Right MI
%
% OBSERVATIONS TO REPORT:
%
% Presence of MRCP:
% - Is there a slow negative shift before cue?
% - What is the amplitude and timing?
%
% Lateralization:
% - Does MRCP differ between left and right MI?
% - Is there contralateral bias?
%
% If MRCP is present:
% - This supports MRCP hypothesis
% - Could explain pre-cue window superiority
%
% If MRCP is absent:
% - Pre-cue activity is likely visual or artifact
%
% STATISTICAL RESULTS:
% - Compare pre-cue amplitude between left and right MI
% - Test if MRCP amplitude differs from zero


% ------------------------------------------------------------
\subsection{Artifact Characterization}
\label{subsec:results_artifacts}

% TARGET: 1 page
%
% CONTENT:
%
% FIGURE: Frontal channel activity
% - Channels: Fp1, Fp2, F7, F8
% - Time-frequency maps or time series
%
% OBSERVATIONS TO REPORT:
%
% Frontal activity:
% - Is there significant activity at frontal sites?
% - Is frontal activity class-discriminative? (different for left vs right)
%
% If lateralized frontal activity present:
% - Strong indication of eye movement artifacts
% - Subjects may track falling apple with eyes
%
% Time-locked patterns:
% - Any transients that could indicate eye movements?
% - When do they occur relative to cue?
%
% QUANTITATIVE:
% - Report t-test comparing Fp1 vs Fp2 between left and right MI
% - If significant: artifact concern


% ============================================================
% 5.2 LEARNED FILTER ANALYSIS (RQ2)
% ============================================================
\section{What EEGNet Learns: Filter Analysis}
\label{sec:results_filters}

% TARGET: 4-5 pages


% ------------------------------------------------------------
\subsection{Temporal Filter Frequency Responses}
\label{subsec:results_temporal_filters}

% TARGET: 2-2.5 pages
%
% CONTENT:
%
% FIGURE: Frequency response of all 8 temporal filters
% - X-axis: Frequency (Hz), 0-40 Hz
% - Y-axis: Magnitude (normalized)
% - 8 subplots or overlaid with different colors
% - Mark mu (8-13) and beta (13-30) bands
%
% OBSERVATIONS TO REPORT:
%
% Filter categorization:
% - Filter 1: Peak at X Hz, category [delta/theta/mu/beta/gamma]
% - Filter 2: Peak at X Hz, category
% - ... (for all 8)
%
% Distribution summary:
% - N filters target mu band: [count]
% - N filters target beta band: [count]
% - N filters target low frequency (<4 Hz): [count]
% - N filters are broadband or other: [count]
%
% INTERPRETATION SETUP (not interpretation itself):
% - If majority target mu/beta: consistent with ERD learning
% - If many target low frequency: consistent with MRCP or artifacts
%
% INCLUDE TABLE:
% - Filter ID | Peak Frequency (Hz) | Category | Bandwidth
%
% ADDITIONAL FIGURE:
% - Heatmap: Filters × Frequency (magnitude as color)


% ------------------------------------------------------------
\subsection{Spatial Filter Topographies}
\label{subsec:results_spatial_filters}

% TARGET: 2-2.5 pages
%
% CONTENT:
%
% FIGURE: Topographic maps for all 16 spatial filters
% - 4×4 grid of topoplots
% - Color scale: filter weights (red positive, blue negative)
% - Consistent color scale across all filters
%
% OBSERVATIONS TO REPORT:
%
% Visual inspection:
% - Which filters focus on sensorimotor region?
% - Which filters focus on frontal region?
% - Which filters focus on occipital region?
% - Which filters are bilateral vs lateralized?
%
% Quantitative metrics:
% - For each filter, compute:
%   * Sensorimotor Ratio (SR)
%   * Frontal Ratio (FR)
%   * Occipital Ratio (OR)
%
% INCLUDE TABLE:
% - Filter ID | SR | FR | OR | Dominant Region
%
% Summary statistics:
% - Mean SR across filters: X ± Y
% - Mean FR across filters: X ± Y
% - Mean OR across filters: X ± Y
%
% ADDITIONAL FIGURE:
% - Bar plot: SR, FR, OR for each filter
% - Or: Scatter plot of SR vs FR (each point is a filter)


% ============================================================
% 5.3 ATTRIBUTION ANALYSIS (RQ3)
% ============================================================
\section{What Drives Classification: Attribution Analysis}
\label{sec:results_attribution}

% TARGET: 5-6 pages


% ------------------------------------------------------------
\subsection{Temporal Attribution Profile}
\label{subsec:results_temporal_attribution}

% TARGET: 1.5-2 pages
%
% CONTENT:
%
% FIGURE: Mean attribution over time
% - X-axis: Time (-1s to 2s relative to cue)
% - Y-axis: Mean |attribution| across channels
% - Separate lines for Left MI and Right MI
% - Mark cue onset (t=0)
% - Show confidence bands (SEM)
%
% OBSERVATIONS TO REPORT:
%
% Temporal distribution:
% - When is attribution highest?
% - Pre-cue (-1 to 0s): X% of total attribution
% - Post-cue (0 to 2s): Y% of total attribution
%
% Pre-cue vs Post-cue:
% - Statistical comparison of attribution magnitude
% - t-test or paired comparison
%
% Peaks and patterns:
% - Any transient peaks? (could indicate artifacts)
% - Sustained attribution? (could indicate ERD)
%
% INCLUDE TABLE:
% - Time Window | Mean Attribution | % of Total


% ------------------------------------------------------------
\subsection{Spatial Attribution Topography}
\label{subsec:results_spatial_attribution}

% TARGET: 2-2.5 pages
%
% CONTENT:
%
% FIGURE: Topographic maps of attribution
% - Panel A: Left MI trials - grand average attribution
% - Panel B: Right MI trials - grand average attribution
% - Panel C: Difference (Left - Right)
%
% OBSERVATIONS TO REPORT:
%
% Spatial distribution:
% - Where is attribution concentrated?
% - Sensorimotor (C3/C4): X% of attribution
% - Frontal (Fp1/Fp2): Y% of attribution
% - Occipital (O1/O2): Z% of attribution
%
% Lateralization:
% - Is Left MI attribution contralateral? (C4 > C3?)
% - Is Right MI attribution contralateral? (C3 > C4?)
%
% STATISTICAL RESULTS:
% - Paired t-test: C3 vs C4 attribution for each class
% - Compare sensorimotor vs frontal attribution
%
% INCLUDE TABLE:
% - Region | Left MI Attribution | Right MI Attribution | Difference
%
% ADDITIONAL FIGURE:
% - Bar plot: attribution by region for each class


% ------------------------------------------------------------
\subsection{Time-Resolved Spatial Attribution}
\label{subsec:results_timeresolved_attribution}

% TARGET: 1-1.5 pages
%
% CONTENT:
%
% FIGURE: Attribution topography at different time points
% - Grid: Rows = time windows, Columns = Left MI / Right MI
% - Time windows: [-1.0,-0.5], [-0.5,0], [0,0.5], [0.5,1.0], [1.0,2.0]
%
% OBSERVATIONS TO REPORT:
%
% Evolution of spatial pattern:
% - Does the spatial pattern change over time?
% - Pre-cue pattern: where is attribution?
% - Post-cue pattern: does it shift to sensorimotor?
%
% This analysis helps distinguish:
% - MRCP: central, early
% - Visual: occipital, early
% - ERD: sensorimotor, sustained
% - Artifact: frontal, early


% ------------------------------------------------------------
\subsection{Comparison to Neurophysiology}
\label{subsec:results_attribution_vs_neuro}

% TARGET: 1 page
%
% CONTENT:
%
% FIGURE: Side-by-side comparison
% - Panel A: ERDS topography (from Section 5.1)
% - Panel B: Attribution topography (from above)
% - Panel C: Spatial correlation plot
%
% QUANTITATIVE COMPARISON:
%
% Spatial correlation:
% - Correlation between ERDS topography and attribution topography
% - r = [value], p = [value]
%
% Interpretation:
% - High correlation: model uses neurophysiological features
% - Low correlation: model may use other features
%
% Channel-by-channel:
% - Scatter plot: ERDS magnitude vs attribution per channel
% - Report correlation


% ============================================================
% 5.4 HYPOTHESIS SCORING
% ============================================================
\section{Evidence Summary: Scoring the Hypotheses}
\label{sec:results_hypothesis_scoring}

% TARGET: 2-3 pages


% ------------------------------------------------------------
\subsection{Summary of Evidence}
\label{subsec:results_evidence_summary}

% TARGET: 1.5-2 pages
%
% CONTENT:
%
% Compile evidence from all previous sections:
%
% INCLUDE TABLE: Hypothesis Scoring Matrix
%
% Rows (Tests):
% - T1: Mu/Beta ERDS at C3/C4
% - T2: MRCP at Cz
% - T3: Frontal activity (artifact marker)
% - T4: Occipital activity (visual marker)
% - T5: Temporal filters target mu/beta
% - T6: Temporal filters target low frequency
% - T7: Spatial filters focus on sensorimotor
% - T8: Spatial filters focus on frontal
% - T9: Attribution at sensorimotor
% - T10: Attribution at frontal
% - T11: Attribution at occipital
% - T12: Attribution highest pre-cue
% - T13: Attribution highest post-cue
%
% Columns: Artifact | ERD | MRCP | Visual
%
% Cells: ++, +, 0, -, --
%
% SCORING LOGIC (briefly explain):
% - T1 ERDS present → supports ERD, against artifact
% - T3 frontal activity → supports artifact
% - etc.


% ------------------------------------------------------------
\subsection{Hypothesis Adjudication}
\label{subsec:results_adjudication}

% TARGET: 0.5-1 page
%
% CONTENT:
%
% Count the evidence:
% - Artifact hypothesis: N tests support, M tests against
% - ERD hypothesis: N tests support, M tests against
% - MRCP hypothesis: N tests support, M tests against
% - Visual hypothesis: N tests support, M tests against
%
% Dominant hypothesis:
% - Which hypothesis has the most support?
% - Is it clearly dominant or is evidence mixed?
%
% NOTE: This is still descriptive. Leave interpretation for Discussion.


% ============================================================
% 5.5 CORRELATION ANALYSIS
% ============================================================
\section{Neural-Performance Correlations}
\label{sec:results_correlations}

% TARGET: 2-3 pages


% ------------------------------------------------------------
\subsection{ERD Magnitude vs Classification Accuracy}
\label{subsec:results_erd_accuracy}

% TARGET: 1-1.5 pages
%
% CONTENT:
%
% FIGURE: Scatter plot
% - X-axis: Mean ERDS magnitude (mu band, C3/C4) per subject
% - Y-axis: Classification accuracy per subject
% - Each point is one subject (N=40)
% - Add regression line
%
% STATISTICAL RESULTS:
% - Pearson r = [value], p = [value]
% - Or Spearman if non-normal
% - 95% CI for correlation
%
% OBSERVATIONS:
% - Is there a significant correlation?
% - If yes: subjects with stronger ERD classified better
% - If no: classification may not rely primarily on ERD


% ------------------------------------------------------------
\subsection{Spatial Alignment vs Accuracy}
\label{subsec:results_alignment_accuracy}

% TARGET: 1 page
%
% CONTENT:
%
% FIGURE: Scatter plot
% - X-axis: Spatial correlation (attribution vs ERDS) per subject
% - Y-axis: Classification accuracy per subject
%
% STATISTICAL RESULTS:
% - Correlation coefficient and p-value
%
% OBSERVATIONS:
% - Does alignment predict performance?


% ------------------------------------------------------------
\subsection{Characterizing BCI Inefficiency}
\label{subsec:results_bci_inefficiency}

% TARGET: 0.5-1 page
%
% CONTENT:
%
% Identify poor performers:
% - Subjects with accuracy < 60% (or other threshold)
% - N = [count] subjects
%
% Compare to good performers:
% - ERDS magnitude: t-test
% - Frontal activity: t-test
% - Other metrics
%
% OBSERVATIONS:
% - What characterizes poor performers?
% - Weaker ERD? More artifacts? Different spatial pattern?


% ============================================================
% 5.6 CHAPTER SUMMARY
% ============================================================
\section{Results Summary}
\label{sec:results_summary}

% TARGET: 0.5 page
%
% CONTENT:
%
% Bullet-point summary of main findings:
% - RQ1: [1-2 sentence summary]
% - RQ2: [1-2 sentence summary]
% - RQ3: [1-2 sentence summary]
% - Dominant hypothesis: [state which]
% - Correlation findings: [brief]
%
% Note: Full interpretation in Discussion chapter
